\section{Finite covers and residual monodromy}
\label{sec:cover}\label{sec:structure}

A smooth semipositive anticanonical metric determines much of the
geometry of the universal cover.  The remaining issue for sections is
the action of the deck group.  We record the geometric reduction with
that action retained, so that an invariant section on the final compact
factor can later be descended to a finite cover of the original variety.
The geometric inputs are the Ricci-semipositive structure theorem of
Demailly--Peternell--Schneider \cite[Structure Theorem and \S3]{DPS1996}
and the rational-connectedness refinement of
Campana--Demailly--Peternell \cite[Theorem~1.4]{CDP2015}.

\begin{theorem}[Finite cover with compact torus monodromy]
\label{thm:structure}
Let $X$ be a smooth connected projective complex variety whose
anticanonical bundle has a smooth Hermitian metric of semipositive
Chern curvature.  There are a connected finite \'etale cover
$\nu:S\to X$, an integer $a\geq0$, and smooth connected projective
varieties $F$ and $Z$ with the following properties.
\begin{enumerate}
\item For a suitable K\"ahler metric on $X$, the common universal cover
of $X$ and $S$ has a holomorphic isometric decomposition
\begin{equation}\label{eq:cover-product}
             \widetilde X=\mathbb C^a\times F\times Z.
\end{equation}
The compact factors are simply connected.  The induced metric on $F$
is Ricci-flat, and $K_F$ has a nowhere-zero parallel holomorphic section.
The anticanonical bundle of $Z$ is smoothly semipositive, and
\begin{equation}\label{eq:cover-noforms}
 H^0(Z,\Omega_Z^p)=0\quad(p>0),\qquad
 H^q(Z,\mathcal O_Z)=0\quad(q>0).
\end{equation}
The variety $Z$ is rationally connected when it has positive dimension.

\item There are a full lattice $\Lambda\subset\mathbb C^a$ and a
homomorphism $\rho:\Lambda\to\operatorname{Aut}(Z)$ for which
\begin{equation}\label{eq:cover-action}
 S=(\mathbb C^a\times F\times Z)/\Lambda,\qquad
 \lambda\cdot(z,f,v)=(z+\lambda,f,\rho(\lambda)v).
\end{equation}
The closure $T$ of $\rho(\Lambda)$ in the holomorphic isometry group
of $Z$ is a compact connected real torus.  The trivial torus is allowed.

\item The action of $T$ extends to an algebraic action of a complex
algebraic torus $G$ on $Z$, with $T$ Zariski dense in $G$.
On $K_Z^{-1}=\det T_Z$ the action is the natural one induced by the
differential.  For every $m\geq0$, its representation on sections satisfies
\begin{equation}\label{eq:cover-invariants}
       H^0(Z,K_Z^{-m})^T=H^0(Z,K_Z^{-m})^G.
\end{equation}

\item The canonical bundles of $\mathbb C^a$ and $F$ have nowhere-zero
holomorphic frames $\eta_0$ and $\eta_F$ invariant under $\Lambda$.
For every integer $m\geq0$, the product projections
$\operatorname{pr}_0$, $\operatorname{pr}_F$, and $\operatorname{pr}_Z$
give the natural, $\Lambda$-equivariant identification
\begin{equation}\label{eq:cover-canonical}
 K_{\widetilde X}^{-m}
 =\operatorname{pr}_0^*K_{\mathbb C^a}^{-m}
   \otimes\operatorname{pr}_F^*K_F^{-m}
   \otimes\operatorname{pr}_Z^*K_Z^{-m},
\end{equation}
in which the first two factors have the invariant frame
$\eta_0^{-m}\otimes\eta_F^{-m}$.
\end{enumerate}
Any of the factors may be a point.  A point contributes the canonical
frame $1$; if $a=0$, the lattice is the zero group.
\end{theorem}

The quotient in \eqref{eq:cover-action} describes a locally trivial
bundle over $\mathbb C^a/\Lambda$.  Its monodromy $\rho$ can have infinite
image, so the finite cover $S$ need not be a product.  This is the reason
for retaining the natural torus action on the anticanonical bundle.

\begin{proof}
We first obtain the product and its cohomological properties.  We then
choose a finite-index deck subgroup and identify its compact action
with an algebraic torus action.

\smallskip\noindent
\emph{The metric splitting and the compact factors.}
The curvature form of the given metric on $K_X^{-1}$ is a smooth
semipositive representative of its first Chern class, with the usual
normalizing factor.  Yau's prescribed-Ricci theorem therefore gives a
K\"ahler metric on $X$ whose Ricci form is that curvature form
\cite{Yau1978}; this application is also stated explicitly in
\cite[p.~218]{DPS1996}.  In particular its Ricci tensor is nonnegative.
The structure theorem of \cite{DPS1996} gives a holomorphic isometric
decomposition
\[
          \widetilde X=\mathbb C^a\times\prod_i M_i,
\]
where each $M_i$ is compact and simply connected and the nonflat
factors are irreducible in the real de Rham decomposition.  The
compactness of these factors is the Cheeger--Gromoll input in the
structure theorem; the holomorphic decomposition comes from the
parallel complex structure.  Group the Ricci-flat factors into $F$
and the remaining factors into $Z$.

Here is the Bochner argument for the part of this decomposition used
below.  On a compact K\"ahler manifold with nonnegative Ricci tensor,
a holomorphic $p$-form $u$ is parallel.  Moreover, if
$\lambda_1,\ldots,\lambda_n\geq0$ are the Ricci eigenvalues in a
unitary frame and $u=\sum_{|J|=p}u_J\,dz_J$, the integrated Bochner
formula forces the pointwise vanishing of
\[
             \sum_{|J|=p}\left(\sum_{j\in J}\lambda_j\right)|u_J|^2.
\]
These are the Bochner consequences used in
\cite[\S2]{DPS1996}; equivalently see the formula in
\cite[(1.6)]{CDP2015}.  In particular, contraction of $u$ with a
positive Ricci eigendirection is zero.  If $u\ne0$ on an irreducible
factor, its contraction kernel
\[
             \{v\in T^{1,0}M_i:\iota_vu=0\}
\]
is a proper parallel complex subbundle.  At any point where the Ricci
tensor is nonzero, that kernel is also nonzero.  Its underlying real
subbundle contradicts irreducibility.  Thus every non-Ricci-flat
factor has no nonzero holomorphic form of positive degree.  The
decomposition of exterior forms on a product and K\"ahler Hodge
symmetry give \eqref{eq:cover-noforms}.  Rational connectedness of
these factors, including the irreducible compact Hermitian-symmetric
case, is the additional conclusion of
\cite[Theorem~1.4(a) and \S4]{CDP2015}; it is not inferred merely from
the vanishing of forms.

The canonical Chern connection of the Ricci-flat factor $F$ is flat.
Since $F$ is simply connected, parallel transport gives a nowhere-zero
parallel holomorphic frame $\eta_F$ of $K_F$.  The metric induced on
$K_Z^{-1}$ has semipositive curvature, since the Ricci tensor of the
product restricts to that of $Z$.

Projectivity of each compact factor follows directly from that of
$X$.  Fix points in the other factors and compose its slice in
$\widetilde X$ with the covering map to $X$.  This holomorphic map
has injective differential.  Pulling back a positive Hermitian metric
on an ample line bundle of $X$ therefore gives a positive line bundle
on the compact factor.  Kodaira's embedding theorem makes that factor
projective \cite{Kodaira1954}.  The argument requires the slice map
to be immersive, not injective.

\smallskip\noindent
\emph{Choosing the deck subgroup.}
Let $\Gamma=\pi_1(X)$ act on $\widetilde X$ by deck transformations.
These are holomorphic isometries.  Uniqueness of the de Rham
decomposition implies that, after a finite-index passage removing
permutations of its factors, every deck transformation is a product
map.  Put $M=F\times Z$.  Projection of this deck subgroup to the
Euclidean isometry group is discrete and has finite kernel.  Indeed,
for a fixed bound on the displacement of $0\in\mathbb C^a$, a
sufficiently large closed Euclidean ball $B$ satisfies
\[
       \gamma(B\times M)\cap(B\times M)\ne\varnothing
\]
for every deck transformation within that bound.  Since $B\times M$
is compact, proper discontinuity allows only finitely many such
transformations.  This proves both assertions.  The Euclidean image
is cocompact, by compactness of the quotient of the full product.
Bieberbach's theorem supplies a finite-index translation lattice in
this image, as in \cite[\S3(ii)]{DPS1996}.

For completeness, the finite kernel does not obstruct choosing an
actual lattice subgroup of the deck group.  Above the translation
lattice we have an extension
\[
       1\longrightarrow N\longrightarrow\Gamma_1
        \longrightarrow\Lambda_1\longrightarrow 0
\]
with $N$ finite.  Pass to the centralizer of $N$ in $\Gamma_1$;
this has finite index because $\operatorname{Aut}(N)$ is finite.
Its kernel $N'$ is central, and its image $\Lambda'$ is a sublattice
of finite index in $\Lambda_1$.  Choose lifts
$g_1,\ldots,g_{2a}$ of a lattice basis $\lambda_1,\ldots,\lambda_{2a}$ of $\Lambda'$.
Their commutators belong to $N'$.  If $e>0$ is divisible by the orders
of all elements of $N'$, then
\[
                  [g_i^e,g_j^e]=[g_i,g_j]^{e^2}=1.
\]
The commuting elements $g_i^e$ generate a free abelian subgroup:
a relation among them projects to a relation among the independent
vectors $e\lambda_i$.  Its image is $e\Lambda'$, so this subgroup
has finite index in the original deck group.  In case $a=0$, the deck
group is finite and we instead take its trivial subgroup.

We next eliminate the action on $F$.  The holomorphic isometry group
of a compact simply connected Ricci-flat manifold is finite.  To see
this, its Lie algebra consists of Killing fields.  The integrated
Bochner identity makes each such field parallel.  Its dual one-form
is closed and, by simple connectedness, exact.  An exact parallel
one-form on a compact manifold is zero: a primitive attains a maximum,
where its differential vanishes.  Thus the compact isometry group has
zero-dimensional Lie algebra and is finite.  Passing to the kernel of
the action on $F$ preserves a finite-index lattice subgroup.

The remaining action on $Z$ has commuting image in its compact
holomorphic isometry group.  Its closure is a compact abelian Lie
group.  Its identity component is a torus and its component group is
finite.  Pass once more to the preimage of that identity component.
The new image is dense in the identity component, since the latter is
open in the old closure.  Denote this connected closure by $T$, the
deck lattice by $\Lambda$, and its action on $Z$ by $\rho$.  This
proves \eqref{eq:cover-action}.

The corresponding quotient $S$ is a connected finite unramified
holomorphic cover of $X$.  The pulled-back positive line bundle
makes $S$ projective, and the covering is a finite \'etale
morphism of projective varieties.  We have now accounted for every
part of the deck action.  In particular $\eta_F$ is invariant because
the action on $F$ is trivial, while
$\eta_0=dz_1\wedge\cdots\wedge dz_a$ is invariant under translations.
The determinant of the product tangent decomposition gives
\eqref{eq:cover-canonical}, with its natural equivariance.

\smallskip\noindent
\emph{The algebraic torus.}
It remains to put the compact action on $Z$ in the algebraic category.
The point case is immediate, so suppose $\dim Z>0$ and choose a very
ample line bundle $A$ on $Z$.  By \eqref{eq:cover-noforms}, the
exponential sequence makes
\[
                 c_1:\operatorname{Pic}(Z)\longrightarrow
                      H^2(Z,\mathbb Z)
\]
injective.  Every element of the connected group $T$ is isotopic to
the identity and hence acts trivially on integral cohomology.
Consequently $t^*A\simeq A$ for every $t\in T$.
The complete linear system of $A$ therefore defines a canonical
projective representation
\[
           T\longrightarrow
           \operatorname{PGL}\bigl(H^0(Z,A)^*\bigr)
\]
preserving the embedded variety $Z$.  Choices of isomorphisms
$t^*A\simeq A$ change the associated linear maps only by scalars,
so this projective representation is well-defined.
It is continuous: choose a projective frame from the nondegenerate
irreducible embedded variety; a projective transformation is
determined continuously by the images of that frame.  Such a frame
exists by first choosing a projective basis and then a point outside
its finitely many coordinate hyperplanes.

Write $r=\dim H^0(Z,A)$.  To diagonalize the compact projective
action, take its compact preimage under the finite covering
$\operatorname{SL}_r(\mathbb C)\to\operatorname{PGL}_r(\mathbb C)$.
The identity component of this preimage is abelian and surjects onto
$T$: its Lie algebra maps isomorphically to the abelian Lie algebra
of $T$, and its image is an open and closed subgroup of $T$.
Averaging a Hermitian form makes this connected preimage unitary;
its commuting matrices can then be diagonalized simultaneously.
Thus, after conjugation, $T$ lies in a diagonal algebraic torus in
$\operatorname{PGL}_r(\mathbb C)$.

Let $G$ be its Zariski closure there.  It is connected, since a
connected $T$ cannot meet distinct components of an algebraic group
without being disconnected.  A connected algebraic subgroup of a
diagonal torus is an algebraic torus.  The projective stabilizer of
$Z$ is Zariski closed, so $G$ preserves $Z$ and acts algebraically.
Differentiation gives the natural algebraic linearization on
$K_Z^{-1}$.  The induced finite-dimensional representation on
$H^0(Z,K_Z^{-m})$ is algebraic.  The stabilizer of any section is
Zariski closed, and Zariski density of $T$ proves
\eqref{eq:cover-invariants}.  All the constructions respect point
factors and the stated canonical-frame conventions.
\end{proof}

\subsection{The norm of a section}

Once a section has been constructed on the finite cover, the following
elementary norm returns it to the original variety.  Its exponent is
the degree of the cover; no uniform degree is needed in the application.

\begin{lemma}[Finite \'etale norm]
\label{lem:etale-norm}
Let $\nu:S\to X$ be a finite \'etale morphism of degree $d>0$
between smooth connected projective complex varieties, and let $A$ be
a line bundle on $X$.  Every nonzero section
$s\in H^0(S,\nu^*A)$ has a nonzero norm
\[
                    \operatorname{Nm}_{\nu}(s)
                    \in H^0(X,A^{\otimes d}).
\]
In particular, a nonzero section of $K_S^{-m}$, for $m>0$, gives a
nonzero section of $K_X^{-md}$.
\end{lemma}

\begin{proof}
Take an analytic open set $U\subset X$ on which $A$ has a frame $e$
and the cover splits into $d$ sheets.  On these sheets write
$s=f_i\nu^*e$ for holomorphic functions $f_1,\ldots,f_d$.
Define the local norm by
\[
                    \left(\prod_{i=1}^d f_i\right)e^{\otimes d}.
\]
This expression is unchanged by a permutation of the sheets.
If the frame is replaced by $e'=u e$, the functions become
$f_i'=f_i/u$; their product changes by $u^{-d}$ and the new tensor
frame changes by $u^d$.  The expressions therefore glue to a
holomorphic section of $A^{\otimes d}$.

Since $S$ is smooth and connected it is integral.  The zero set of the
nonzero section $s$ is thus proper and has dimension less than
$\dim S$ (or is empty).  Its image under the finite map is a proper
closed subset of $X$.  Away from that image every sheetwise factor is
nonzero, so the norm is nonzero.  Projectivity identifies this
holomorphic section with an algebraic section.  Finally,
$K_S=\nu^*K_X$ because $\nu$ is \'etale; applying the result
with $A=K_X^{-m}$ proves the last assertion.  In dimension zero the
same construction is the product of nonzero constants.
\end{proof}
